\section{Examples of Generated Verifiable Tasks}
\label{app:generated-task-examples}

This section shows representative tasks generated by \ours across different
desktop applications. Each task is paired with executable verification criteria
that check the resulting application state, files, metadata, or persistent side
effects.

\begin{taskexample}{Zotero: Three-Level Collection Hierarchy}
\taskfield{Application} \texttt{zotero}

\taskfield{Task} Launch Zotero with the pre-seeded library at
\texttt{/home/user/Zotero/zotero.sqlite}. Under \texttt{My Library}, create a
top-level collection named \texttt{Papers}, a subcollection named
\texttt{Vision}, and a sub-subcollection named \texttt{Object Detection}.
Each collection must be a direct child of the previous level. Leave Zotero
running so the updated library state is saved.

\taskfield{Initial files}
\texttt{zotero.sqlite} at \texttt{/home/user/Zotero/zotero.sqlite}.

\taskfield{Representative verification checks}
\begin{compactitem}
    \item The collection \texttt{Papers} exists in the Zotero library.
    \item The collection \texttt{Vision} exists.
    \item The collection \texttt{Object Detection} exists.
    \item \texttt{Vision} is a direct child of \texttt{Papers}.
    \item \texttt{Object Detection} is a direct child of \texttt{Vision}.
\end{compactitem}
\end{taskexample}


\begin{taskexample}{MuseScore: Piano Solo to Piano Quartet}
\taskfield{Application} \texttt{musescore3}

\taskfield{Task} Open the score
\texttt{/home/user/Documents/piano\_sketch.mscz}, which contains a single
Piano part in C major, 4/4 time, and four measures of quarter notes. Arrange
the solo as a piano quartet by adding \texttt{Violin}, \texttt{Viola}, and
\texttt{Violoncello} parts below the Piano. Copy the Piano melody to the
Violin staff. On the Viola staff, enter one whole note per measure with pitches
\texttt{A3}, \texttt{D4}, \texttt{G3}, and \texttt{C4}. On the Violoncello
staff, enter a walking bass using four quarter notes per measure:
\texttt{C3}, \texttt{E3}, \texttt{G3}, \texttt{E3};
\texttt{F3}, \texttt{A3}, \texttt{C4}, \texttt{A3};
\texttt{G3}, \texttt{B3}, \texttt{D4}, \texttt{B3};
and \texttt{C3}, \texttt{E3}, \texttt{G3}, \texttt{C4}. Add dynamics
\texttt{mp}, \texttt{mf}, \texttt{p}, and \texttt{f} to measures 1--4,
respectively. Add staccato articulations to all four Violoncello notes in
measure 4 and nowhere else. Save the score, then export it as uncompressed
MusicXML to \texttt{/home/user/Documents/piano\_sketch.musicxml}.

\taskfield{Initial files}
\texttt{piano\_sketch.mscz} at
\texttt{/home/user/Documents/piano\_sketch.mscz}.

\taskfield{Representative verification checks}
\begin{compactitem}
    \item The original \texttt{.mscz} file still exists at the expected path.
    \item The saved score has exactly four parts: \texttt{Piano},
    \texttt{Violin}, \texttt{Viola}, and \texttt{Violoncello}.
    \item The score still has four measures and remains in 4/4 time.
    \item The score contains at least 52 notes, covering the original Piano
    melody, copied Violin melody, Viola whole notes, and Violoncello bass line.
    \item Measures 1--4 contain the expected dynamics:
    \texttt{mp}, \texttt{mf}, \texttt{p}, and \texttt{f}.
    \item Measure 4 contains staccato articulations.
    \item The total number of staccato articulations in the score is exactly 4.
    \item The exported MusicXML file exists at
    \texttt{/home/user/Documents/piano\_sketch.musicxml}.
    \item The exported MusicXML file contains exactly four parts.
\end{compactitem}
\end{taskexample}


\begin{taskexample}{LibreOffice Calc: Sales Commission Spreadsheet}
\taskfield{Application} \texttt{libreoffice\_calc}

\taskfield{Task} Open \texttt{/home/user/Documents/commissions.xlsx}. In the
\texttt{Sales} sheet, add bold headers \texttt{Commission Rate} and
\texttt{Commission Amount}. Use a nested \texttt{IF} formula to assign
commission rates based on monthly sales: 10\% for sales above 20,000, 8\% for
sales at least 10,000, and 5\% otherwise. Compute commission amounts for all
20 records, create a new \texttt{Commission Summary} sheet, and add formulas
for total sales, total commission, and average commission rate. Save the
workbook.

\taskfield{Initial files}
\texttt{commissions.xlsx} at \texttt{/home/user/Documents/commissions.xlsx}.

\taskfield{Representative verification checks}
\begin{compactitem}
    \item Cells \texttt{D1} and \texttt{E1} contain the expected bold headers.
    \item Cell \texttt{D2} contains an \texttt{IF} formula and evaluates to
    the expected commission rate.
    \item Cell \texttt{E2} references \texttt{C2} and \texttt{D2} and computes
    the expected commission amount.
    \item The copied formulas produce expected values for low-, mid-, and
    high-tier sales examples.
    \item The \texttt{Commission Summary} sheet exists.
    \item The summary sheet contains the expected labels and formulas for total
    sales, total commission, and average commission rate.
    \item The workbook is saved.
\end{compactitem}
\end{taskexample}


\begin{taskexample}{Obsidian: Recipe Vault Construction}
\taskfield{Application} \texttt{obsidian}

\taskfield{Task} Open the Obsidian vault at
\texttt{/home/user/Documents/RecipeVault}. Create folders for
\texttt{Italian}, \texttt{Asian}, and \texttt{Desserts}. Add four recipe notes
with YAML frontmatter, headings, ingredients, instructions, tags, and internal
links: \texttt{Carbonara.md}, \texttt{Cacio e Pepe.md},
\texttt{PadThai.md}, and \texttt{Tiramisu.md}. Finally, create a root
\texttt{Index.md} note linking to all recipes and tagged as an index.

\taskfield{Initial files}
None.

\taskfield{Representative verification checks}
\begin{compactitem}
    \item The \texttt{Italian}, \texttt{Asian}, and \texttt{Desserts} folders
    exist in the vault.
    \item Recipe notes contain the expected frontmatter fields, such as
    \texttt{cuisine}, \texttt{servings}, and \texttt{time}.
    \item \texttt{Carbonara} links to \texttt{Cacio e Pepe}, and
    \texttt{Cacio e Pepe} links back to \texttt{Carbonara}.
    \item \texttt{PadThai} contains the \texttt{\#noodles} tag.
    \item \texttt{Tiramisu} links to \texttt{Carbonara}.
    \item \texttt{Index.md} links to the expected recipes and contains the
    \texttt{\#index} tag.
    \item The vault contains exactly five notes.
\end{compactitem}
\end{taskexample}


\begin{taskexample}{Blender: Parent-Child Rig Hierarchy}
\taskfield{Application} \texttt{blender}

\taskfield{Task} Open the Blender file
\texttt{/home/user/Documents/rig\_hierarchy.blend}, which starts empty. Add an
Empty and rename it to \texttt{Rig}. Add a cube named \texttt{Torso}, a UV
sphere named \texttt{Head}, and a cylinder named \texttt{Arm}. Parent
\texttt{Torso} and \texttt{Arm} to \texttt{Rig}, and parent \texttt{Head} to
\texttt{Torso}, making \texttt{Head} a grandchild of \texttt{Rig}. Save the
file.

\taskfield{Initial files}
\texttt{rig\_hierarchy.blend} at
\texttt{/home/user/Documents/rig\_hierarchy.blend}.

\taskfield{Representative verification checks}
\begin{compactitem}
    \item The saved Blender file exists at the expected path.
    \item The object \texttt{Rig} exists and has type \texttt{EMPTY}.
    \item The objects \texttt{Torso}, \texttt{Head}, and \texttt{Arm} exist and
    have mesh object types.
    \item \texttt{Torso} is parented to \texttt{Rig}.
    \item \texttt{Arm} is parented to \texttt{Rig}.
    \item \texttt{Head} is parented to \texttt{Torso}.
\end{compactitem}
\end{taskexample}